\documentclass[12pt]{article}
\usepackage[utf8]{inputenc}
\usepackage[T1]{fontenc}
\usepackage{hyperref}
\title{Papavô Language Resources for K-12 Curriculum Development}
\author{Compiled by Claude Sonnet 5.5 (claude-sonnet-5-5) for Andr\'as Kornai}
\date{}
\begin{document}
\maketitle
\section{Introduction}
Papavô (ISO 639-3 ppv) is an unattested language of Brazil. No linguistic documentation, word list, speaker community or educational material was found in the searches, so no curriculum can be based on it.
\section{Primary Educational Resources}
None found.
\section{Linguistic and Reference Materials}
None found.
\section{Religious and Community Resources}
\begin{itemize}
\item Joshua Project lists a Papavo people group in Brazil with unknown or very small population, possibly merged into neighbouring groups (a missionary-oriented source, not a linguistic one): \url{https://joshuaproject.net/people_groups/14306/BR}
\end{itemize}
\section{Digital and Audio Resources}
None found.
\section{Curriculum Development Context}
With no attested data and no identified community, a K-12 curriculum is not feasible. This dossier is intentionally short.
\end{document}